\documentclass[12pt,a4paper]{article}
\usepackage[T2A]{fontenc}
\usepackage[utf8]{inputenc}
\usepackage[russian]{babel}
\usepackage{geometry}
\usepackage{hyperref}
\usepackage{graphicx}
\usepackage{listings}
\usepackage{xcolor}
\usepackage{booktabs}
\geometry{left=2.5cm,right=2.5cm,top=2.5cm,bottom=2.5cm}

\lstset{
  basicstyle=\ttfamily\small,
  keywordstyle=\color{blue},
  stringstyle=\color{red},
  commentstyle=\color{gray},
  breaklines=true,
  showstringspaces=false
}

\begin{document}

\begin{center}
\textbf{САНКТ-ПЕТЕРБУРГСКИЙ НАЦИОНАЛЬНЫЙ}\\
\textbf{ИССЛЕДОВАТЕЛЬСКИЙ УНИВЕРСИТЕТ ИТМО}\\[1.5cm]

Дисциплина: Фронт-энд разработка\\[1.5cm]

Отчет\\[1.5cm]

Домашняя работа 5\\[2cm]

Выполнил:\\
Черников Кирилл\\
Группа J3210\\[0.8cm]

Проверил:\\
Добряков Д. И.\\[2cm]

Санкт-Петербург\\
2026 г.
\end{center}

\newpage

\section{Цель работы}

Изучить основные команды менеджера зависимостей NPM и освоить базовые концепции фреймворка Vue.js~3: создание проекта через Vite, компонентный подход, реактивность, передача данных через пропсы и пользовательские события.

\section{Требования задания}

\begin{itemize}
  \item Изучить основные команды NPM (\texttt{npm init}, \texttt{npm install}, \texttt{npm run})
  \item Создать Vue-проект с помощью Vite
  \item Реализовать несколько компонентов, демонстрирующих возможности Vue.js
  \item Предоставить скриншоты компонентов и код приложения
\end{itemize}

\section{Ход выполнения работы}

\subsection{Основные команды NPM}

NPM (Node Package Manager) — стандартный менеджер пакетов для Node.js. Основные команды, использованные в работе:

\begin{center}
\begin{tabular}{ll}
\toprule
Команда & Описание \\
\midrule
\texttt{npm init}              & Инициализация нового проекта, создание \texttt{package.json} \\
\texttt{npm install}           & Установка всех зависимостей из \texttt{package.json} \\
\texttt{npm install \textless пакет\textgreater}   & Установка конкретного пакета в \texttt{dependencies} \\
\texttt{npm install -D \textless пакет\textgreater} & Установка пакета в \texttt{devDependencies} \\
\texttt{npm run dev}           & Запуск dev-сервера (скрипт из \texttt{package.json}) \\
\texttt{npm run build}         & Сборка проекта для продакшена \\
\bottomrule
\end{tabular}
\end{center}

\subsection{Создание проекта}

Проект создан вручную с использованием Vite и плагина \texttt{@vitejs/plugin-vue}. Файл \texttt{package.json}:

\begin{lstlisting}[language={}]
{
  "name": "dataforge-hw5",
  "scripts": {
    "dev":     "vite",
    "build":   "vite build",
    "preview": "vite preview"
  },
  "dependencies":    { "vue": "^3.4.27" },
  "devDependencies": {
    "@vitejs/plugin-vue": "^5.0.4",
    "vite": "^5.2.11"
  }
}
\end{lstlisting}

Конфигурация Vite (\texttt{vite.config.js}):

\begin{lstlisting}[language={}]
import { defineConfig } from 'vite'
import vue from '@vitejs/plugin-vue'

export default defineConfig({
  plugins: [vue()],
})
\end{lstlisting}

Структура проекта:

\begin{lstlisting}
vue-app/
  index.html
  vite.config.js
  package.json
  src/
    main.js
    App.vue
    assets/styles.css
    components/
      UserGreeting.vue
      CounterWidget.vue
      TaskList.vue
\end{lstlisting}

Точка входа \texttt{src/main.js}:

\begin{lstlisting}[language={}]
import { createApp } from 'vue'
import App from './App.vue'
import './assets/styles.css'

createApp(App).mount('#app')
\end{lstlisting}

\subsection{Компонент UserGreeting — пропсы и события}

Компонент принимает имя и роль через пропсы и отправляет событие \texttt{update:role} при смене роли через выпадающий список. Это демонстрирует однонаправленный поток данных и паттерн \texttt{v-model} на пользовательских компонентах.

\begin{lstlisting}[language={}]
<template>
  <div class="card">
    <h3>Привет, <strong>{{ name }}</strong>!</h3>
    <div class="role-tag">{{ role }}</div>
    <select :value="role"
            @change="$emit('update:role', $event.target.value)">
      <option>Исполнитель</option>
      <option>Валидатор</option>
      <option>Менеджер</option>
    </select>
  </div>
</template>

<script setup>
defineProps({ name: String, role: String })
defineEmits(['update:role'])
</script>
\end{lstlisting}

\subsection{Компонент CounterWidget — реактивность}

Компонент демонстрирует реактивность Vue через \texttt{ref}. Значение счётчика меняется при нажатии на кнопки, интерфейс обновляется автоматически.

\begin{lstlisting}[language={}]
<template>
  <div class="card">
    <div class="counter-value">{{ count }}</div>
    <div class="btn-row">
      <button @click="count++">+ Добавить</button>
      <button :disabled="count === 0" @click="count--">
        − Убрать
      </button>
      <button :disabled="count === 0" @click="count = 0">
        Сброс
      </button>
    </div>
  </div>
</template>

<script setup>
import { ref } from 'vue'
const count = ref(0)
</script>
\end{lstlisting}

\subsection{Компонент TaskList — директивы v-for и v-if}

Компонент реализует список задач и демонстрирует:
\begin{itemize}
  \item \texttt{v-for} — рендеринг списка задач
  \item \texttt{v-if} / \texttt{v-else} — условный рендеринг пустого состояния
  \item \texttt{v-model} — двустороннее связывание поля ввода и чекбоксов
  \item \texttt{computed} — автоматическое вычисление числа выполненных задач
\end{itemize}

\begin{lstlisting}[language={}]
<template>
  <div class="card">
    <form @submit.prevent="addTask">
      <input v-model="newTask" placeholder="Новая задача..." />
      <button type="submit">Добавить</button>
    </form>

    <p v-if="tasks.length === 0">Список пуст.</p>

    <ul v-else>
      <li v-for="task in tasks" :key="task.id"
          :class="{ done: task.done }">
        <input type="checkbox" v-model="task.done" />
        <label>{{ task.text }}</label>
        <button @click="removeTask(task.id)">X</button>
      </li>
    </ul>
    <p>Выполнено: {{ doneCount }} / {{ tasks.length }}</p>
  </div>
</template>

<script setup>
import { ref, computed } from 'vue'

const newTask = ref('')
const tasks = ref([
  { id: 1, text: 'Изучить основы Vue.js', done: true },
  { id: 2, text: 'Создать первый компонент', done: false },
])
let nextId = 3

const doneCount = computed(
  () => tasks.value.filter(t => t.done).length
)

function addTask() {
  const text = newTask.value.trim()
  if (!text) return
  tasks.value.push({ id: nextId++, text, done: false })
  newTask.value = ''
}

function removeTask(id) {
  tasks.value = tasks.value.filter(t => t.id !== id)
}
</script>
\end{lstlisting}

\subsection{Корневой компонент App.vue}

Корневой компонент импортирует три дочерних компонента и передаёт данные через пропсы. Смена роли в \texttt{UserGreeting} обновляет состояние в родителе через \texttt{v-model}.

\begin{lstlisting}[language={}]
<template>
  <div class="app">
    <header class="app-header">
      <h1>DataForge — изучение Vue.js</h1>
    </header>
    <div class="grid">
      <UserGreeting :name="userName" v-model:role="userRole" />
      <CounterWidget />
    </div>
    <TaskList />
  </div>
</template>

<script setup>
import { ref } from 'vue'
import UserGreeting from './components/UserGreeting.vue'
import CounterWidget from './components/CounterWidget.vue'
import TaskList from './components/TaskList.vue'

const userName = ref('Кирилл')
const userRole = ref('Исполнитель')
</script>
\end{lstlisting}

\section{Скриншоты}

\begin{figure}[h!]
  \centering
  \includegraphics[width=0.85\textwidth]{"screenshots/Скриншот 11-05-2026 234858.jpg"}
  \caption{Общий вид приложения: компоненты UserGreeting (пропсы, смена роли), CounterWidget (реактивность) и TaskList (v-for, v-if) на одной странице}
\end{figure}

\begin{figure}[h!]
  \centering
  \includegraphics[width=0.6\textwidth]{"screenshots/Скриншот 11-05-2026 234929.jpg"}
  \caption{Компонент CounterWidget — счётчик показывает значение 5 после пяти нажатий кнопки «+ Добавить»; кнопки «− Убрать» и «Сброс» становятся активными}
\end{figure}

\begin{figure}[h!]
  \centering
  
\includegraphics[width=0.75\textwidth]{"screenshots/Скриншот 11-05-2026 234942.jpg"}
  \caption{Компонент TaskList — задача «Изучить основы Vue.js» отмечена как выполненная (зачёркнута через v-if/v-bind:class), счётчик «Выполнено: 1~/ 2» обновляется реактивно через computed}
\end{figure}

\section{Вывод}

В ходе выполнения домашней работы изучены основные команды NPM и базовые концепции Vue.js~3. Создан Vue-проект на базе Vite с тремя компонентами, демонстрирующими: передачу данных через пропсы (\texttt{UserGreeting}), реактивность через \texttt{ref} (\texttt{CounterWidget}), условный рендеринг и рендеринг списков через \texttt{v-if} и \texttt{v-for} (\texttt{TaskList}). Все концепции, изученные в ДЗ5, применены в полноценном SPA в рамках ЛР3.

\end{document}
